\documentclass[10pt]{article}
\usepackage[margin=0.72in]{geometry}
\usepackage{amsmath,amssymb,amsthm,booktabs,microtype}
\usepackage[hidelinks]{hyperref}
\hypersetup{
  pdftitle={Two-Page Theorem Summary: Stationary Diffusion-Driven Instability in Mass-Action Networks},
  pdfauthor={Alec Kriebel},
  pdfsubject={External audit summary},
  pdfkeywords={mass-action kinetics, stationary Turing instability, complexity}
}

\pagestyle{plain}
\newcommand{\diag}{\operatorname{diag}}
\newcommand{\im}{\operatorname{im}}
\newcommand{\R}{\mathbb R}
\title{Two-Page Theorem Summary\\\large Stationary Diffusion-Driven Instability in Classical Mass Action}
\author{External-audit version}
\date{13 August 2026}
\begin{document}
\maketitle
\vspace{-1.5em}

\paragraph{Decision problem.}
The input is a rationally encoded finite indexed reaction network
$y_r\to y'_r$, with source matrix $Y=[y_r]$ and stoichiometric matrix
$\Gamma=[y'_r-y_r]$.  Write $S=\im\Gamma$.  Decide whether there are strictly
positive rate constants, a strictly positive homogeneous equilibrium, and
strictly positive diagonal diffusion coefficients for all species such that
the homogeneous Jacobian is Hurwitz on $S$ and some nonzero Neumann mode has a
strictly positive real eigenvalue.

\paragraph{Exact characterization.}
For $v>0$, $\Gamma v=0$, and $h>0$, set
\[
 A(v)=\Gamma\diag(v)Y^T,
 \qquad J(v,h)=A(v)\diag(h).
\]
Choose any rational basis $B$ of $S$ and rational left inverse $C$, and set
$R(v,h)=CJ(v,h)B$.  The answer is YES exactly when there are $v,h$ and a
nonempty principal set $I$ such that
\[
 \Gamma v=0,
 \qquad R(v,h)\text{ is Hurwitz},
 \qquad (-1)^{|I|}\det A(v)[I,I]<0.
\]
The converse mass-action realization is
$x_i^*=h_i^{-1}$ and $k_r=v_r/(x^*)^{y_r}$.

The fixed-matrix step is
\[
 \exists D\succ0:\ J-D\text{ has a positive real eigenvalue}
 \iff
 \exists I:\ (-1)^{|I|}\det J[I,I]<0.
\]
It follows from the principal-minor expansion of $\det(D-J)$ and a multiscale
assignment that makes one negative coefficient uniquely dominant.  Column
scaling gives
\[
 (-1)^{|I|}\det J[I,I]
 =\Bigl(\prod_{i\in I}h_i\Bigr)(-1)^{|I|}\det A(v)[I,I].
\]

\paragraph{Conservation.}
If $J(\R^n)\subseteq S$ and $J|_S$ is Hurwitz, then
$\R^n=S\oplus\ker J$ and the conservation zero eigenvalues are semisimple.
A nonzero Neumann eigenfunction has zero spatial mean, so its species
amplitude is unrestricted; one must test $J-q^2D$ on the full species space.

\newpage
\paragraph{Complexity.}
A Boolean selector $b_i^2-b_i=0$, with $S_b=\diag(b)$, satisfies
\[
 \det(I-S_b-S_bA(v)S_b)=(-1)^{|I|}\det A(v)[I,I],
 \qquad I=\{i:b_i=1\}.
\]
Determinant and Hurwitz circuits have polynomial size, so the decision problem
belongs to the existential theory of the reals.

The problem is NP-hard under a polynomial-time many-one reduction from
\textsc{Partition}.  For $m=k^2>\ell$, the source family is
\[
 Q=I+aa^T,
 \quad \beta=1-\frac1{2m(1+a^Ta)},
 \quad B(x,y)=\begin{pmatrix}-kQ&y\\x^T&k\beta\end{pmatrix}.
\]
Its open cube contains a Hurwitz matrix exactly for YES partition instances.
A row-splitting lift $L(q)$ is similar to
$\left(\begin{smallmatrix}B(q)&U\\0&-I\end{smallmatrix}\right)$.  A separate
determinant argument proves
\[
 \exists q,D\succ0:\ L(q)D\text{ Hurwitz}
 \iff \textsc{Partition}=\text{YES},
\]
so arbitrary equilibrium/right scaling cannot create a false YES.  An exact
row-segment reaction construction realizes the full family
$\{R L(q):R\succ0\}$ with independent steady-flux constraints.  The network
has $3m+1$ species, $8m+2$ reactions, full stoichiometric rank, and columns
$\pm e_i$.  Its complex molecularity is unbounded, and the hardness is weak.

\paragraph{Fixed species count.}
Let $K=\{v\geq0:\Gamma v=0\}$ and
$P=\{\Gamma\diag(v)Y^T:v\in K\}$.  If a positive flux exists, positive
fluxes project exactly to $\operatorname{ri}P$.  Every extreme ray of $K$ is a
positive circuit supported on at most $\operatorname{rank}\Gamma+1\leq n+1$
reactions.  For fixed $n$ these rays, the fixed-dimensional projected cone,
and a relative-facet description are computable in polynomial time.  The
remaining real formula has a fixed number of variables, giving polynomial bit
complexity for every fixed $n$.

\paragraph{Certificates and limitations.}
YES instances have finite real-algebraic sample certificates; NO instances
have finite Real-Nullstellensatz identities.  No polynomial size or practical
general generator is claimed.  The theorem concerns weak stationary linear
instability only: it does not require a first, simple, or transverse crossing,
does not exclude an earlier wave instability, does not prove a nonlinear
patterned state, and does not give a finite motif characterization.
\end{document}
